\clearpage
\fancyhead[L]{\small OBSERVATIONS ASTRONOMIQUES.}
\hypertarget{native-646-lower}{}
\begin{center}\small Göttingische gelehrte Anzeigen. Stück 129. Seite 1281..1284. 1826 August 14.\\[3pt]\rule{.72\linewidth}{.3pt}\end{center}

\emph{Observations astronomiques faites à l'observatoire royal de Paris, publiées par le bureau des longitudes.} Paris 1825. Bei Bachelier. 402 Seiten in Folio.

Die Beobachtungen, welche auf der Pariser Sternwarte in den Jahren 1800..1809 angestellt sind, sind in den Bänden der Connaissance des tems 1808..1812 und 1823..1825 bekannt gemacht; die spätern sollen nach einem Beschluss des Bureau des longitudes besonders gedruckt werden, und der vorliegende erste Band enthält, mit Ausschluss der Beobachtungen am dreifussigen \textsc{Reichenbach}'schen Repetitionskreise, diejenigen, die in den Jahren 1810..1819 von vier Beobachtern, den Hrn. \textsc{Bouvard}, \textsc{Arago}, \textsc{Mathieu} und \textsc{Nicollet} angestellt sind. Die zu den Beobachtungen angewandten Instrumente werden in der Einleitung beschrieben. Das Mittagsfernrohr, angefangen von \textsc{Ramsden} und vollendet von \textsc{Berge}, wurde im August 1803 aufgestellt. Es hat $7\frac{1}{2}$ Fuss Brennweite und 4 Zoll Oeffnung; die Länge der Axe ist 4 Fuss; die Vergrösserung nicht ganz eine hundertmalige. Seitenbeweglichkeit des Oculars und Beleuchtung durch die Axe, wie jetzt allgemein gewöhnlich ist. Zur Berichtigung des Mittagsfernrohrs sind zwei Meridianzeichen errichtet, das nördliche in der Entfernung von 1364 Meter auf dem Palais du Luxembourg, das südliche an einer Pyramide in der Ebne von Montrouge, 1840 Meter entfernt. Zu Zielpunkten dienen kreisrunde Löcher in Metallplatten, die in horizontaler Richtung verschiebbar sind, von $2\frac{1}{4}$ Zoll $\cdot$ Durchmesser; das nördliche projicirt sich gegen eine dahinter gestellte weissgefärbte Platte von Eisenblech, das südliche gegen den Himmel; die Durchmesser erscheinen $9''07$ und $7''1$ gross. Verticalfäden sind fünf, deren Zwischenräume als genau gleich angesehen und von Sternen im Aequator in $17^{\rm s}36$ durchlaufen werden; aus welcher Materie sie bestehen, wird nicht erwähnt. Die Beobachtungen an diesem Instrumente nehmen mehr als die Hälfte des Bandes ein; ihre Gegenstände sind die Sonne, der Mond, und die meisten Planeten und die \textsc{Maskelyne}'schen Fixsterne, selten andere; von den Cometen von 1811 und 1819 kommen einige untere Culminationen vor. Was übrigens die Beobachtungen selbst betrifft, so werden die Forderungen, die die neuere Astronomie an selbstständige Beobachtungen mit Instrumenten von so ausgezeichneten Dimensionen macht, nicht ganz befriedigt; sie bieten keine zureichende Mittel dar, zur Untersuchung, mit welchem Grade von Genauigkeit die Meridianzeichen in der Mittagsfläche sich befinden; bei der geringen Entfernung dieser Zeichen und der grossen Brennweite des Fernrohrs müssen sie mit grosser Parallaxe gegen die Fäden und geringer Deutlichkeit erscheinen. .Dazu kommt, dass sie im Ganzen nicht oft zur Prüfung angewandt sind, z. B. im Jahre 1810, nachdem in den frühern Monaten öfters bemerkt ist, dass der Nebel gehindert habe sie zu sehen, zum ersten Male den 11. Mai. Wenn die ungünstige Localität die Benutzung der Meridianzeichen so selten verstattete, so wäre eine häufige Beobachtung von Circumpolarsternen doppelt nothwendig gewesen; allein nur selten ist einmal der Polarstern, und noch seltener sind aufeinanderfolgende Culminationen \SourceLeaf{647}oberhalb und unterhalb des Pols beobachtet. Die Horizontalität der Axe ist zwar öfters geprüft; allein Mittel zur Prüfung, ob die beiden Zapfen gleiche Dicke haben, was man selbst bei den Instrumenten von den ersten Künstlern nicht voraussetzen darf, fehlen gänzlich. Auch die Rechtwinklichkeit der optischen Axe zur Drehungsaxe ist selten geprüft. Beobachtungen von Sternbedeckungen, die in diesem Tagebuche der Beobachtungen am Mittagsfernrohr eingeschaltet sind, kommen in grosser Anzahl vor; man muss um so mehr bedauern, dass die angeführten Umstände immer einige kleine Ungewissheit in der Bestimmung der absoluten Zeit zurücklassen, da die Astronomen gewohnt sind, die Angabe der geographischen Länge immer auf die Pariser Sternwarte zu beziehen.

\begingroup\rightskip=0pt plus 1em Den zweiten Abschnitt des Werks machen die Beobachtungen an $7\frac{1}{2}$fussigen \textsc{Bird}'schen Mauerquadranten aus. Das Fernrohr ist gleichfalls $7\frac{1}{2}$ Fuss lang, hat $2\frac{1}{4}$ Zoll Oeffnung und vergrössert 70..80 Mal. Die Beobachtungsgegenstände sind dieselben, wie am Mittagsfernrohr. Den Collimationsfehler hat man für eine Anzahl von Punkten aus zahlreichen Beobachtungen von Fixsternen zu bestimmen gesucht, indem man für deren Declinationen ein Mittel aus den Angaben mehrerer Astronomen zum Grunde legte. Es ist überflüssig, den Rang anzudeuten, welchen die Beobachtungen mit diesem Instrumente bei dem gegenwärtigen Zustande der praktischen Astronomie haben. Dieser Quadrant dient für die südliche Hälfte des Meridians; an der Nordseite ist ein zweiter fünffussiger von \textsc{Sisson} aufgehängt, derselbe, welchen einst \textsc{Lalande} in Berlin gebrauchte. Der vorliegende Band enthält jedoch keine Beobachtungen mit diesem Instrumente. --- Seit dem Jahr 1823 ist ein Mauerkreis von \textsc{Fortin} im Gebrauch, dessen Beschreibung wir im nächsten Bande zu erwarten haben.\par\endgroup

Der dritte Abschnitt enthält die Beobachtungen an der parallactischen Maschine von \textsc{Bellet}, welche in einem abgesonderten Theile der Sternwarte aufgestellt ist. Die Durchmesser der beiden Kreise derselben sind 13 Zoll. Das Fernrohr hat drei Fuss Länge, beinahe $2\frac{1}{2}$ Zoll Brennweite und vergrössert 40..50 Mal. Die Beobachtungen betreffen theils die in dem Zeitraum von 1810..1819 erschienenen Cometen, theils den Mondsflecken Manilius, behuf einer Untersuchung der Libration des Mondes, deren sehr schätzbare Resultate bekanntlich schon vor mehreren Jahren in der Connaissance des tems bekannt gemacht sind. Dieses Instrument ist im Jahre 1823 an die Sternwarte in Marseille abgegeben, und statt desselben ein grosses Aequatoreal von \textsc{Gambey} aufgestellt, dessen Beschreibung im nächsten Bande folgen soll.
